\section{Reference contract and execution model}
\label{sec:contract}
\subsection{Output qualification}
Let $\rho$ identify the reference model and tokenizer bytes, prompt construction, numerical runtime, ordinary execution recipe, sampler and stopping rule. For input $x$, its endpoint is
\begin{equation}
\mathcal O_\rho(x)=(\mathbf y,\mathbf e,n,y_{\mathrm{last}},s),
\label{eq:endpoint}
\end{equation}
where $\mathbf y$ is the full native token-ID array, $\mathbf e$ its end-of-generation (EOG) flags, $n$ its length and $s$ the stop reason. The principal campaigns stop at native EOG or 128 generated tokens. A candidate qualifies on a registered panel $X$ when every required pair has $\mathcal O_c(x)=\mathcal O_\rho(x)$, the reference is stable and all required records are present. A shorter divergent response cannot qualify a latency comparison.

We distinguish exact-arithmetic equivalence, bitwise operator correspondence and endpoint agreement. Each answers a different question. Backend acceptance and actual operator placement are additional execution prerequisites. These definitions apply throughout the paper; qualification always refers to the stated configuration, panel and output window.

\subsection{The stateful transaction}
A target context owns KV values, sequence positions, a pending sampled token or available logits, and relevant sampler state. The draft has separate state. It proposes up to $k$ tokens; the target evaluates the pending token and proposal positions; acceptance, correction and any bonus determine the committed prefix. Rejected positions are reconciled before the next cycle. A visible token may still be pending evaluation into KV and must be consumed exactly once.

QB changes the target attention representation inside this transaction. It adds no numerical-consistency verifier; the ordinary speculative verification and reconciliation still execute. Microbatch capacity $U$, occupied query count $q$, caller width and output reservation remain distinct. We write UB$n$ for a configured capacity of $n$. A caller containing $k+1$ rows can execute as several physical micrographs, so caller counts alone do not describe parallel work.

\subsection{Local repair and state composition}
The design objective is to minimize complete request time over tested implementations, subject to endpoint qualification. Write a request fragment as $D_g(P_g(x))$, where $P_g$ forms the state consumed by $D_g$ under execution organization $g$. Correspondence of $D_{\QB}$ and the reference on one input state $s_\rho$ does not establish correspondence on the different state $P_g(x)$.

For full stateful composition, a relation between candidate and reference states must hold initially, preserve the committed output at each transaction, and hold again after reconciliation. This is a standard sufficient-condition argument; its premises are tested empirically here. Section~\ref{sec:state} examines two simpler assumptions often used in local reasoning: that endpoint-equivalent states have interchangeable components, and that enlarging a sufficient replacement preserves its effect.
